% Generated draft: requires mathematical and source review.
\documentclass[11pt]{article}
\usepackage[margin=25mm]{geometry}
\usepackage{fontspec}
\setmainfont{Latin Modern Roman}
\usepackage{amsmath,amssymb,mathtools,unicode-math}
\setmathfont{Latin Modern Math}
\usepackage{enumitem,needspace,xurl,hyperref}
\setlength{\parindent}{0pt}
\setlength{\parskip}{5pt}
\setlength{\emergencystretch}{3em}
\newcommand{\N}{\mathbb{N}}
\newcommand{\Z}{\mathbb{Z}}
\newcommand{\Q}{\mathbb{Q}}
\newcommand{\R}{\mathbb{R}}
\newcommand{\C}{\mathbb{C}}
\newcommand{\F}{\mathbb{F}}
\newcommand{\A}{\mathbb{A}}
\newcommand{\PP}{\mathbb{P}}
\newcommand{\E}{\mathbb{E}}
\newcommand{\eps}{\varepsilon}
\newcommand{\dd}{\,\mathrm{d}}
\begin{document}
\section*{UnsolvedMath ID: 20001994}
\textbf{Source problem code:} catalogue entry
\subsection{Definitions and mathematical statement}
A nontrivial commutative idempotent semifield is a commutative semiring $(K,\oplus,\odot,0,1)$ with $0\ne1$, $a\oplus a=a$, and every nonzero element multiplicatively invertible. Its natural order $a\le b$ if $a\oplus b=b$ need not be total. The source's ``algebraically closed'' terminology requires an idempotent notion of algebraic equation/root; it is not interpreted as ordinary zero roots of all nonconstant polynomials in a zero-sum-free semiring.

\textbf{A.12, six-part program of G.~L.~Litvinov in cooperation with G.~B.~Shpiz.} The main statement gives only part 1; parts 2--6 occur in the nested remark and are part of this record:
\begin{enumerate}
\item Construct algebraic geometry over a class of algebraically closed idempotent semifields broader than tropical semifields.
\item Define abstract algebraic varieties, beyond only affine or projective ones, over tropical and idempotent semifields.
\item Define regular functions and maps giving a natural category of affine varieties and a corresponding, contravariant coordinate-semiring description in the spirit of classical algebraic geometry.
\item Define equations and ideal analogues so that points and subvarieties correspond to appropriate algebraic data.
\item Define prime objects and irreducible varieties and investigate decomposition into irreducible components.
\item Construct dequantization procedures relating classical and tropical varieties and seek a functor, or the source's less precise ``almost functor'', between suitable categories.
\end{enumerate}
These are requests for compatible foundations, not a single fully axiomatized conjecture. The categories, root convention and meaning of almost functor are not fixed. The remark attributes an algebraic-closedness criterion to Shpiz (2000) and says that some standard linear function spaces and all Banach lattices generate examples, referring to Litvinov--Maslov--Shpiz (2001); the generating construction and criterion are not supplied here.

\textbf{Supplied local obstruction and congruence repair.} For arbitrary such $K$, without algebraic closedness or a total order, use the formal polynomial semiring $K[x_1,\ldots,x_n]$ and a point $a\in K^n$. Define
\[
 I_0(a)=\{f:f(a)=0\},\qquad
 \ker(\operatorname{ev}_a)=\{(f,g):f(a)=g(a)\}.
\]
A semiring congruence is an equivalence relation compatible with addition and multiplication; its kernel is a relation on pairs, not an ordinary ideal. The aid asserts
\[
 I_0(a)=\langle x_i:a_i=0\rangle,\qquad
 \ker(\operatorname{ev}_a)=\langle x_i\sim a_i:1\le i\le n\rangle_{\mathrm{cong}}.
\]
In particular $I_0(a)=\{0\}$ for every torus point $a\in(K\setminus\{0\})^n$, so ordinary zero-vanishing ideals retain only which coordinates vanish. By contrast the evaluation congruence determines $a$ injectively because constants are preserved.

In one variable, for $c\in K\setminus\{0\}$, the bend congruence of $x\oplus c$ is generated by deleting each of its two monomials:
\[
 x\oplus c\sim x,\qquad x\oplus c\sim c.
\]
The aid asserts
\[
 \operatorname{Bend}(x\oplus c)=\langle x\sim c\rangle_{\mathrm{cong}}
                          =\ker(\operatorname{ev}_c).
\]
This exact point-level comparison addresses an obstruction to an ideal-only implementation of part 4. It does not provide abstract gluing, a general Nullstellensatz, a prime-spectrum theory, decomposition theorems or a functorial dequantization for all the source's proposed semifields. Formal polynomials are distinguished from polynomial functions, and no particular modern foundational framework is silently chosen as the unique solution to the six-part program.
\subsection{Short English statement}
Litvinov and Shpiz ask for a broad algebraic geometry over idempotent semifields, including varieties, maps, prime objects and dequantization. The supplied aid isolates a point-level issue: ordinary zero ideals forget nonzero coordinates, while evaluation and bend congruences recover them. This local repair does not settle the entire foundational program.
\subsection{Sources}
\begin{itemize}
\item Local UnsolvedMath registry, permanent ID 20001994, source code catalogue entry; canonical \texttt{aim-geometry-notes.json}, notes[331], including its full nested remark.
\item AIM \emph{Amoebas and tropical geometry}, A.12 ``Idempotent geometry'', contributed by G.~L.~Litvinov in cooperation with G.~B.~Shpiz: \url{https://aimath.org/WWN/amoebas/amoebas.pdf}.
\item G.~Shpiz, ``Solving algebraic equations in idempotent semifields'', \emph{Uspekhi Mat. Nauk}, vol.~55, no.~5 (2000), pp.~185--186, in Russian; the source notes an English translation in \emph{Russian Mathematical Surveys} (2000).
\item G.~L.~Litvinov, V.~P.~Maslov and G.~B.~Shpiz, ``Idempotent functional analysis: an algebraic approach'', \emph{Math. Notes}, vol.~69, no.~5 (2001), pp.~758--797, bibliographic data as supplied, not independently corrected.
\item Registry research attempt 1, dated 2026-08-14, partial result, rigor ``proved'', novelty confidence low. The machine-generated, non-peer-reviewed aid requires independent verification before citation.
\end{itemize}
\end{document}
